\subsection*{A larger certified gain using pairs of primes}

For the integer problem a refinement gives the absolute lower bound
\[
                         d(E)-d(G)>10^{-728}.
\]
Use the same base-nine split and pivot prime $2$. Write $D=\{0,1,3,4\}$
for the old low digits. The spokes are the $265$ odd primes through
$1699$. Every new low vector has pivot $2$, with spoke digits $1$ or
$2$. Include the vector with all spoke digits $1$, all vectors with
one spoke digit $2$, and some vectors with two spoke digits $2$. The
selection of pairs is specified below.

For an old low vector, call a spoke active if its digit differs from
$1$. Delete precisely the following old vectors, with $a\in D$ the
pivot digit.
\begin{itemize}
\item With no active spoke, delete the vector if $a\in\{1,3,4\}$.
\item With one active spoke, delete its vector if the spoke digit
is $4$, or if it is $3$ and $a\ne0$.
\item With two active spokes $p<r$, delete their vector if $\{p,r\}$
is selected and their ordered digits belong to
$\{3,4\}\times\{3,4\}$ or equal $(0,4)$. In this line the selection
condition applies to both alternatives.
\end{itemize}
Keep every remaining old vector.

\paragraph{A finite proof of avoidance.}
Two new vectors in a putative progression force the third pivot to be
$2$, and the cube $\{1,2\}^{266}$ contains only constant progressions.
If there is just one new vector, every spoke where it has digit $1$ must
also be $1$ at both old points, since $D$ has no nonconstant arithmetic
progression. It suffices to check the pivot and the $k=0,1,2$ spokes
where the new vector equals $2$. For $k=2$ that pair qualifies for the
two-spoke deletion rule because its new vector was selected. Write $R$
for the test in the three bullets on these ordered coordinates, with
the pair considered selected. For each $v\in D^{k+1}$ form
\[
                  u=4\mathbf{1}-v
             \quad\text{and}\quad
                  u=2v-2\mathbf{1}.
\]
Retain $\{u,v\}$ when $u\in D^{k+1}$. The first choice treats
$2\mathbf{1}$ as the middle, the second as an endpoint with old
middle $v$. The numbers of distinct unordered pairs are respectively
$4,12,40$ for $k=0,1,2$. In every pair $R(u)$ or $R(v)$ holds.
This last finite identity is verified by the loops in
\texttt{check\_spoke\_pairs.py}. At most $4^3$ values of $v$ occur in
one case. This exhausts all signed progressions with one new point.
All-old progressions were already excluded. The base-nine lifting
preserves this conclusion since each digit sum and each doubled digit
is at most eight.

\paragraph{Selecting and pricing each pair.}
Divide the low-digit weights by the weight of the all-ones vector.
Summing over the four possible old pivot digits gives the multiplier
$2+1+1/4+1/8=27/8$. For $p<r$ the net contribution of selecting the
pair is
\[
 J(p,r)=\frac1{2pr}
 -\frac{27}{8}\left[
       (p^{-2}+p^{-3})(r^{-2}+r^{-3})+\frac{p}{r^3}\right].
\]
The last term inside the brackets is the weight of digits $(0,4)$.
Select exactly the pairs with $J(p,r)>0$. Their additions and
deletions have distinct active spoke sets, so their contributions add
without overlaps. The complete normalized gain is
\[
 B=-\frac78+
 \sum_{3\le p\le1699}\left(\frac1{2p}
                   -\frac{11}{8p^2}-\frac{27}{8p^3}\right)
             +\sum_{3\le p<r\le1699}\max\{0,J(p,r)\},
\]
where the indices in both sums are primes. Here $11/8$ is the sum of
the normalized weights of pivot digits $1,3,4$, which are the three
pivot values deleted when the single active spoke has digit $3$.

The integer checker finds $15{,}693$ selected pairs. It bounds $B$ from
below by rounding every summand down to a multiple of $10^{-15}$ and
obtains
\[
                          B\ge\frac{73099512201}{10^{15}}.
\]
For independent reproduction,
$8p^3r^3 J(p,r)=4p^2r^2-27((p+1)(r+1)+p^4)$. Put $P=\prod_{p\le1699}p$. The low-block weight
normalization including the geometric prefactors is $\varphi(P)/P^2$.
The independent higher-digit weights contribute
$\prod_{p\mid P}\sum_{h\in\Ag}p^{-9h}\ge1$. The outside greedy
density is at least $1/\zeta(2)>1/2$. Finally the checker verifies
the integer comparison
\[
       73099512201\,\varphi(P)\,10^{728}>2\cdot10^{15} P^2.
\]
Multiplying these bounds gives the claimed improvement in natural
density. No floating-point arithmetic enters the certificate.

\paragraph{The actual size of the exchange.}
Let $\eta=(d(E)-d(G))/d(G)$. The identity
$H(w)=(1+w)(1+w^3)H(w^9)$ cancels both the outside factors and all high-digit factors when dividing \eqref{eq:delta} by $d(G)$. Thus the \emph{relative} increase is an exact rational number,
\[
 \eta=B\prod_{p\le1699}\frac{p^3}{(p+1)(p^3+1)} .
\]
Computing $B$ from the displayed finite sum as a rational gives
$2.185320021572<10^{727}\eta<2.185320021573$.
We also certify
\begin{equation}\label{eq:numericgain}
              1.5728 <10^{727}\bigl(d(E)-d(G)\bigr)< 1.5730 .
\end{equation}
For the absolute interval, take $T=100000$ and multiply the factors
$(1-p^{-1})(1+p^{-1})(1+p^{-3})(1+p^{-9})$ for $p\le T$.
This retains the first three ternary positions. Relative to this truncated probability, the true $d(G)$ can be lower by at most
$\sum_{n>T}n^{-2}\le1/T$, or higher by at most
$\sum_{n\ge2}n^{-27}\le 2^{-27}+2^{-26}/26$.
The checker bounds the truncated probability by rounding each product outwards to $18$ decimal places. It tests both ends of \eqref{eq:numericgain} by rational comparisons. This completes Theorem~\ref{th:main}.

For an example of the actual rule, let $P$ be the primorial through $1699$. The exchange drops $P$ and admits $2P$, $2pP$ for each spoke, and $2prP$ for each selected pair. The high digits and primes outside the finite list supply all the other integers represented by these patterns. The change in membership therefore has a genuine limiting density even though the numerical gain is exceptionally small.
